%% Numbers in this section come from generated/numbers.tex and
%% generated/table_*.tex (built by tools/make_assets.py from the fetched
%% experiment artifacts). Do not type results here by hand.
\section{Evaluation}
\label{eval}

We evaluate \textsc{lgm} on three axes that mirror its intended role in
TFOW: (i)~one-step predictive quality of the next event's type and timing,
against neural and non-parametric baselines
(Sec.~\ref{sec:eval-set}); (ii)~closed-loop
stability of the learned dynamics and generative realism of free-running
simulations under the stylized facts of high-frequency order flow
(Sec.~\ref{sec:eval-stylized}); and (iii)~price-level realism of
simulated books against recorded ground truth
(Sec.~\ref{sec:eval-price}).  Together these establish the model as a
faithful generative model of LOB microstructure---the property the 3S
detection and query-localization stages rely on.

We model each arrival as a \emph{single} event type, i.e.\ the
singleton restriction $|x|{=}1$ of the set-valued process of
Sec.~\ref{subsec:mark_head_choice}; predictions and likelihoods are over
the next event type, not over subsets.

\subsection{Evaluation Setup}\label{sec:eval-setup}

\paragraph{Data.}
We use Level-2 order book and trade data from two centralized crypto
exchanges, \textbf{Gemini} and \textbf{Coinbase} (Kaiko feed),
reconstructed into atomic events by the event constructor of Section~3
(MO/CO/LO/IS $\times$ side $\times$ top $K{=}10$ levels; 62 fixed event
types).  Sliding windows of 50 events (stride 32) are split
chronologically \emph{within each (token, venue) pair} 70/15/15 into
train/validation/test, with a one-window gap at each boundary to
prevent leakage.  Timestamps are exchange-stamped on a $0.1$\,s grid.
Each arrival is reduced to a single event type (the singleton case of
Sec.~\ref{subsec:mark_head_choice}).

\paragraph{Models.}
Our proposed model is \textbf{\textsc{lgm}} (linear rate-pinned
multi-timescale Hawkes ground rate $\times$ deep soft-max marks,
Sec.~\ref{subsec:hawkes_mean_calibration}).  We compare against seven
baselines under an identical protocol (shared single-item type head,
multi-hot event histories, and the same training budget---30 epochs,
batch 512, AdamW, $10^{-3}$): two classic marked point processes, a
parametric \textsc{Hawkes} process~\cite{Hawkes1971} and
\textsc{RMTPP}~\cite{du2016recurrent}; a structured state-space decoder
\textsc{s2p2}; and four neural baselines following the borrowed-backbone
protocol of Shi and Cartlidge~\cite{Shi2022StateSimulation}---an
\textsc{lstm}, a self-attentive Transformer
(\textsc{sahp}~\cite{zhang2020self}), a continuous-time decayed-cell
network (\textsc{ct-lstm}~\cite{Mei2017TheProcess}), and a per-type
parallel recurrent network
(\textsc{pct-lstm}~\cite{Shi2022StateSimulation}).\footnote{The neural
baselines are \emph{adapted} backbones under a shared single-item
type-decoding protocol, not paper-faithful re-implementations;
comparisons probe the architectures' inductive biases, not the original
systems.}  A non-parametric empirical-transition baseline
(\textsc{EmpTrans}; next-type frequencies conditioned on the last event
type) anchors the tables.

\paragraph{Protocol.}
All headline metrics are \emph{target-free}: no model receives the
ground-truth next-event time.  The predictive distribution over the
next arrival time is obtained from the learned ground intensity by
trapezoid quadrature of the compensator
$\Lambda(t)=\int_0^{t}\lambda(s)\,ds$ on a log-spaced grid; point
predictions read out the decision-theoretically matched statistic
(predictive median for MAE, predictive mean for RMSE).  The next event
\emph{type} is scored from the conditional mark distribution
$p^*(k\mid t)$ evaluated at the model-predicted time: we report top-1
type accuracy and per-event type perplexity (the head-agnostic
single-item likelihood, comparable across all models).

\subsection{Next-Event Prediction}\label{sec:eval-set}

Table~\ref{tab:predict} compares single-item next-event prediction
across all models: next-type top-1 accuracy and perplexity, and
target-free next-time MAE/RMSE.  All models share the single-item type
head and histories, so the comparison isolates the effect of the
intensity/backbone.  \emph{Results are pending the Gemini/Coinbase
experiment runs} (Sec.~\ref{sec:eval-setup}); the table records the
model set and metrics.

\begin{table}[t]
\caption{Single-item next-event prediction on the held-out test set:
next-type top-1 accuracy ($\uparrow$), type perplexity ($\downarrow$),
and target-free next-time MAE/RMSE ($\downarrow$). \textsc{lgm} is the
proposed model; the remainder are baselines. Best in \textbf{bold}.
Cells marked ``--'' are pending experiment runs.}
\label{tab:predict}
\small
\setlength{\tabcolsep}{4pt}
\begin{tabular}{lcccc}
\toprule
Model & Acc $\uparrow$ & Perplexity $\downarrow$ & Time MAE $\downarrow$ & Time RMSE $\downarrow$ \\
\midrule
EmpTrans (non-param.)        & -- & -- & -- & -- \\
Hawkes~\cite{Hawkes1971}     & -- & -- & -- & -- \\
RMTPP~\cite{du2016recurrent} & -- & -- & -- & -- \\
LSTM                         & -- & -- & -- & -- \\
SAHP~\cite{zhang2020self}    & -- & -- & -- & -- \\
CT-LSTM~\cite{Mei2017TheProcess} & -- & -- & -- & -- \\
s2p2                         & -- & -- & -- & -- \\
PCT-LSTM~\cite{Shi2022StateSimulation} & -- & -- & -- & -- \\
\midrule
\textbf{LGM} (ours)          & -- & -- & -- & -- \\
\bottomrule
\end{tabular}
\end{table}

\subsection{Stylized Facts of Simulated Order Flow}\label{sec:eval-stylized}

This and the price-realism section evaluate \textsc{lgm} against the
baselines in closed loop; the numbers below reflect the earlier
diagnostic decoders and will be regenerated with \textsc{lgm} as the
proposed model on the Gemini/Coinbase runs. \textsc{lgm}'s rate pin
(exact stationary mean, Sec.~\ref{subsec:hawkes_mean_calibration}) is
precisely the property expected to deliver closed-loop stability, the
criterion this section isolates.

A generative model is only a \emph{simulator} if it remains
on-distribution when its own outputs are fed back as history.  We
evaluate this in closed loop: event sets, inter-arrival times, and
per-channel volumes are sampled from the model (times by inversion of
the compensator, sets from the conditional set distribution, volumes
from the log-volume head) and appended to the conditioning history,
for a fixed simulated \emph{duration} (600\,s) so that decoders with
different event rates are compared over equal market time.  Rollouts
neither explode nor die out: simulated inter-arrival and set-size
marginals remain at a constant distance from the real marginals across
the full horizon.

Following the realism-evaluation practice of market
simulators~\cite{vyetrenko2020getreal}, we aggregate real and simulated
streams into 1-second buckets and compute the eleven stylized facts of
Cont~\cite{cont2001empirical} on a signed order-flow proxy for returns
(price-moving events: market orders and best-level in-spread
insertions) and an activity series (events per bucket).
Figure~\ref{fig:stylized} shows the panel for \textsc{vsm-rmtpp}.
The decoders separate sharply.  \textsc{vsm-rmtpp} is the best
all-round simulator: near-zero return autocorrelation
(\sfRmtppAcfAbsSim{} vs.\ \sfRmtppAcfAbsReal{} real), the closest
tails (kurtosis \sfRmtppSkewReal-level), correct skew sign, a matching
activity--volatility coupling (\sfRmtppActVolSim{} vs.\
\sfRmtppActVolReal{} real), and the only non-degenerate volatility
clustering signal (ACF$|r|$ \sfRmtppAcfAbsSim{} vs.\
\sfRmtppAcfAbsReal{} real).  \textsc{vsm-hawkes} under joint volume
training simulates almost no volatility memory
(\sfHawkesAcfAbsSim); an ablation without the volume head restores
its clustering to $0.051$ and its kurtosis to near-real ($1.89$ vs.\
$1.93$), isolating a multi-task interference effect: the volume
likelihood's gradients, flowing through the shared recurrent cell,
dampen the \emph{variability} of the self-exciting intensity while
leaving prediction metrics unchanged.  \textsc{vsm-s2p2} is
degenerate in closed loop (return autocorrelation $0.89$: a
near-deterministic drift), examined next.

\begin{figure}[t]
\centering
\autofig{figs/stylized_facts_rmtpp.png}{\linewidth}
\caption{Cont's eleven stylized facts, real vs.\ simulated order flow
for \textsc{vsm-rmtpp} (single seed, 600\,s closed-loop rollouts).}
\label{fig:stylized}
\end{figure}

\subsection{Price-Level Realism}\label{sec:eval-price}

\paragraph{From events to prices.}
Simulated streams are converted to prices by a deterministic
book-replay engine that inverts the event constructor's semantics:
market orders consume the touch and walk the ladder; cancellations
that empty the best level move the quote away; in-spread insertions
improve it by their tick distance.  Events within one timestamp are
applied in matching-engine order (MO, IS, CO, LO).  The engine assumes
a dense one-tick ladder at the top of book; following the
tick-size-regime analysis of Jain et al.~\cite{Jain2025NoBooks}, this
is accurate for large-tick regimes and an approximation for small-tick
assets, so we convert ladder units to prices via recorded per-level
gaps and restrict price-level claims to distributional statements.
The engine contains no stochastic or fitted components: unlike prior
work that attaches order prices and volumes to simulated events from
separately fitted power laws with hand-tuned aggression
percentages~\cite{Shi2022StateSimulation}, every stochastic
quantity---type, side, book level, aggression, timing, volume---is
produced by the model, and only deterministic mechanics are imposed.

\paragraph{Recorded ground truth.}
The real side of every comparison uses the \emph{recorded} book state
emitted by our event constructor (mid-price and full ten-level ladders
per event), not a reconstruction.  Simulated books are initialized
from the recorded snapshot at the seed point.  Model events that are
inconsistent with the simulated book (cancellations of empty levels,
insertions into a one-tick spread) are skipped and \emph{counted}; we
report this invalid-event rate as a measure of the model's structural
coherence (\textsc{vsm-hawkes} \pvHawkesInvalid\%,
\textsc{vsm-rmtpp} \pvRmtppInvalid\%, \textsc{vsm-s2p2}
\pvSspInvalid\%).

\paragraph{Results.}
Against recorded reality (mean spread \pvHawkesSpreadRealMean\,bps),
the simulated books separate sharply: \textsc{vsm-rmtpp} is the most
faithful (\pvRmtppSpreadSimMean\,bps, with zero impossible in-spread
insertions), \textsc{vsm-hawkes} is over-tight and over-calm
(\pvHawkesSpreadSimMean\,bps), and \textsc{vsm-s2p2}'s book
\emph{diverges} to \pvSspSpreadSimMean\,bps mean
(\pvSspSpreadSimPnf\,bps at p95) --- $40\times$ reality.  Thus the
decoder that wins every one-step prediction column produces the least
viable market: one-step predictive quality does \emph{not} predict
closed-loop viability, a compounding-error phenomenon familiar from
model-based reinforcement learning that we propose as a first-class
evaluation criterion --- \emph{closed-loop book stability} --- for
LOB generative models.

\begin{figure}[t]
\centering
\autofig{figs/price_v2_hawkes.png}{\linewidth}
\caption{Price-level realism for \textsc{vsm-hawkes} against recorded
book state: mid trajectories, return Q-Q, autocorrelations, volatility
clustering, spread, signature plot, book shape, and touch-queue
density.}
\label{fig:pricev2}
\end{figure}

\subsection{Discussion}\label{sec:eval-discussion}

Four findings shape how TFOW should be used.  (1)~\emph{Set-valued
modeling is the source of predictive advantage}: the state-space
variant beats every adapted baseline on set metrics and matches the
frequency table on single-type accuracy --- multi-label structure,
not next-type guessing, is where learnable signal lives.
(2)~\emph{No decoder dominates; choose by task}: the state-space
decoder for prediction and likelihood scoring, the constant-hazard
recurrent decoder for timing and simulation --- its rigid hazard is
simultaneously why it cannot win the prediction table and why it is
the only decoder whose closed-loop market remains faithful.
(3)~\emph{Joint volume modeling is prediction-neutral but dampens
self-exciting dynamics}: volume accuracy matches dedicated baseline
heads with no cost to set/type/timing metrics, yet its gradients,
flowing through the shared recurrent state, suppress the Hawkes
decoder's closed-loop volatility clustering (restored by a
volume-head-free ablation); we report this multi-task interference as
a finding, and a stop-gradient volume readout as the remedy.
(4)~\emph{Known limitations}: the feed's $0.1$\,s grid makes
continuous hazards mis-specified (the cause of both temporal PIT
over-dispersion and the expressive decoders' between-event intensity
inflation; a grid-aware discrete-time hazard is the principled fix),
the conditionally independent set head caps within-tick co-occurrence
structure, and the study covers a single venue, one week, and three
symbols; scaling the corpus, a discrete-time temporal head, and a
within-tick dependent set head are the subject of ongoing work.
